\documentclass[../../../include/open-logic-section]{subfiles}

\begin{document}

\olfileid[id]{sth}{story}{blv}

\olsection{Lampiran: Hukum Dasar V Frege}

Dalam \olref[rus]{sec}, kita menjelaskan bahwa Russell merumuskan
paradoksnya sebagai masalah bagi sistem yang diuraikan Frege dalam
\emph{Grundgesetze}. Sistem Frege tidak memuat perumusan langsung Komprehensi
Naif. Jadi, dalam lampiran ini, kita akan menjelaskan dengan sangat singkat
apa yang \emph{termuat} dalam sistem Frege, bagaimana hubungannya dengan
Komprehensi Naif, dan bagaimana hubungannya dengan Paradoks Russell.

Sistem Frege bersifat \emph{orde kedua} dan dirancang untuk merumuskan
gagasan tentang \emph{ekstensi suatu konsep}.\footnote{Tepatnya, Frege
berusaha memformalkan suatu gagasan yang lebih umum: ``rentang-nilai'' suatu
fungsi. Ekstensi konsep merupakan kasus khusus dari gagasan yang lebih umum
tersebut. Lihat \citet[pp.\ 8--9]{Heck2012} untuk perinciannya.} Dengan
menggunakan notasi yang diilhami Frege, kita akan menulis
$\fregeext{x}{F(x)}$ untuk \emph{ekstensi konsep~$F$}. Perangkat ini mengambil
sebuah \emph{predikat}, ``$F$'', dan mengubahnya menjadi sebuah
\emph{term} (orde pertama), ``$\fregeext{x}{F(x)}$''. Dengan perangkat ini,
Frege memberikan \emph{definisi} keanggotaan berikut:
\[
	a \in b =_\text{df} \exists G(b = \fregeext{x}{G(x)} \land Ga)
\]
secara kasar: $a \in b$ jika dan hanya jika $a$ berada di bawah suatu konsep
yang ekstensinya adalah $b$. (Perhatikan bahwa penguantifikasi ``$\exists G$'' bersifat orde
kedua.) Frege juga mempertahankan prinsip berikut, yang dikenal sebagai
\emph{Hukum Dasar V}:
$$\fregeext{x}{F(x)} = \fregeext{x}{G(x)} \liff \forall x (Fx \liff Gx)$$
secara kasar: konsep-konsep mempunyai ekstensi identik jika dan hanya jika
konsep-konsep itu koekstensif. (Sekali lagi, baik ``$F$'' maupun ``$G$'' berada pada posisi
predikat.) Kini suatu prinsip sederhana menghubungkan keanggotaan dengan
terpenuhinya sifat:

\begin{lem}[dalam \emph{Grundgesetze}]\ollabel{lem:Fregeextensions}
$\forall F \forall a(a \in \fregeext{x}{F(x)} \liff Fa)$
\end{lem}

\begin{proof}
Tetapkan $F$ dan $a$. Sekarang $a \in \fregeext{x}{F(x)}$ jika dan hanya jika
$\exists G(\fregeext{x}{F(x)} = \fregeext{x}{G(x)} \land Ga)$ (menurut
definisi keanggotaan) jika dan hanya jika
$\exists G(\forall x(Fx \liff Gx) \land Ga)$ (menurut Hukum Dasar V) jika
dan hanya jika
$Fa$ (menurut logika orde kedua elementer).
\end{proof}

Dan hasil ini memberikan Komprehensi Naif hampir seketika:

\begin{lem}[dalam \emph{Grundgesetze}.]
$\forall F \exists s \forall a (a \in s \liff Fa)$
\end{lem}

\begin{proof}
Tetapkan $F$; sekarang \olref{lem:Fregeextensions} memberikan
$\forall a (a \in \fregeext{x}{F(x)} \liff Fa)$; jadi
$\exists s\forall a(a \in s \liff Fa)$ berdasarkan generalisasi
eksistensial. Hasilnya mengikuti karena $F$ dipilih secara sembarang.
\end{proof}

Paradoks Russell mengikuti dengan mengambil $F$ sebagaimana diberikan oleh
$\forall x(Fx \liff x \notin x)$.

\end{document}
